% ═══════════════════════════════════════════════════════════════
% LEHRERHINWEISE: Halbwertszeit (Physik Klasse 10)
% Stunde 5-6 (90 Minuten)
% ═══════════════════════════════════════════════════════════════

\documentclass[11pt, a4paper]{article}

\usepackage[utf8]{inputenc}
\usepackage[T1]{fontenc}
\usepackage[ngerman]{babel}
\usepackage{helvet}
\renewcommand{\familydefault}{\sfdefault}

\usepackage[margin=2cm]{geometry}
\usepackage{enumitem}
\usepackage{tabularx}
\usepackage{booktabs}
\usepackage{array}
\usepackage{xcolor}
\usepackage{tcolorbox}
\usepackage{fancyhdr}

\setlist{nosep, leftmargin=1.5em}

% Farben
\definecolor{physik}{HTML}{2c5aa0}
\definecolor{grau}{HTML}{888888}
\definecolor{hellgrau}{HTML}{f0f0f0}
\definecolor{warnung}{HTML}{b35c00}
\definecolor{warnungbg}{HTML}{fff3e0}

% Kopfzeile
\pagestyle{fancy}
\fancyhf{}
\renewcommand{\headrulewidth}{0.4pt}
\lhead{\small\textbf{LEHRERHINWEISE} -- Physik Klasse 10}
\rhead{\small Halbwertszeit (Stunde 5-6)}
\cfoot{\small\thepage}

% Boxen
\newtcolorbox{warnbox}[1][]{
    colback=warnungbg, colframe=warnung, boxrule=0.8pt, arc=0pt,
    left=6pt, right=6pt, top=6pt, bottom=6pt,
    fonttitle=\bfseries\small, title=#1
}

\newcommand{\abschnitt}[1]{%
    \vspace{10pt}\noindent\textbf{\textcolor{physik}{\large #1}}\vspace{6pt}\hrule\vspace{8pt}
}

\begin{document}

\begin{center}
{\LARGE\bfseries Lehrerhinweise: Halbwertszeit}\\[6pt]
{\large Physik Klasse 10 -- Stunde 5-6 (90 Min)}\\[4pt]
{\small\textcolor{grau}{Kernphysik-Einheit, Teil 3}}
\end{center}
\vspace{8pt}\hrule\vspace{12pt}

% ═══════════════════════════════════════════════════════════════
\abschnitt{1. Stundenverlauf (90 Minuten)}

\begin{tabularx}{\textwidth}{|c|c|X|l|}
\hline
\textbf{Zeit} & \textbf{Phase} & \textbf{Inhalt} & \textbf{Material} \\
\hline
0-5 & Einstieg & Frage: Wie lange dauert es, bis Strahlung weg ist? & Tafel \\
\hline
5-15 & Experiment & Würfelzerfall-Simulation (Klasse wirft Würfel) & 30+ Würfel \\
\hline
15-25 & Auswertung & Ergebnis graphisch darstellen → Zerfallskurve & Tafel/Plakat \\
\hline
25-35 & Input & Lernpfad: Definition HWZ, Zerfallsgesetz & LP, Beamer \\
\hline
35-45 & Erarbeitung & AB Kasten 1-2 ausfüllen & AB-03 \\
\hline
45-55 & Input & Lernpfad: Beispiel-HWZ, C-14-Methode & LP \\
\hline
55-65 & Erarbeitung & AB Kasten 3, Übung 1 (Tabelle) & AB-03 \\
\hline
65-75 & Simulation & PhET: Radioactive Dating Game & Tablets/PCs \\
\hline
75-85 & Vertiefung & Übungen 2-4, C-14-Anwendung & AB-03 \\
\hline
85-90 & Sicherung & Zusammenfassung: Warum ist HWZ so wichtig? & Tafel \\
\hline
\end{tabularx}

\vspace{12pt}

% ═══════════════════════════════════════════════════════════════
\abschnitt{2. Würfelzerfall-Experiment}

\textbf{Durchführung:}
\begin{enumerate}
\item Jeder Schüler erhält 1-2 Würfel (mind. 30 insgesamt)
\item Alle würfeln gleichzeitig. Wer eine 6 würfelt: Würfel „zerfallen" → rauslegen
\item Übrige Würfel zählen, notieren, nächste Runde
\item Nach 10-15 Runden: Ergebnisse in Diagramm eintragen
\end{enumerate}

\textbf{Erwartung:} Bei 36 Würfeln nach jeder Runde ca. 1/6 weniger → entspricht nicht genau HWZ, aber zeigt exponentiellen Abfall. HWZ ≈ 3,8 Runden (bei 1/6-Wahrscheinlichkeit).

\vspace{12pt}

% ═══════════════════════════════════════════════════════════════
\abschnitt{3. Differenzierung}

\begin{tabularx}{\textwidth}{|c|X|X|}
\hline
\textbf{Niveau} & \textbf{Fokus} & \textbf{Hilfestellung} \\
\hline
★ (BR) & Tabelle ausfüllen (halbieren), qualitativ & Halbierungskette visualisieren \\
\hline
★★ (MR) & + Berechnungen mit ganzen HWZ & Formelkarte, Taschenrechner \\
\hline
★★★ (GY) & + C-14-Methode, Transfer & Logarithmen optional \\
\hline
\end{tabularx}

\vspace{12pt}

% ═══════════════════════════════════════════════════════════════
\abschnitt{4. Häufige Fehler \& Reaktionen}

\begin{warnbox}[Typische Schülerfehler]
\begin{enumerate}
\item \textbf{„Nach 2 Halbwertszeiten ist alles weg"}\\
→ Korrektur: Nach 2 HWZ sind noch 25\% da! Die Kurve erreicht nie 0.

\item \textbf{Lineare statt exponentielle Abnahme}\\
→ Visualisierung: 100→50→25→12,5 (halbieren!), nicht 100→50→0

\item \textbf{HWZ mit Lebensdauer verwechselt}\\
→ Klarstellung: HWZ ist die Zeit für 50\% Zerfall, nicht für kompletten Zerfall.

\item \textbf{„Schnelle HWZ = gefährlicher"}\\
→ Differenzieren: Kurze HWZ = hohe Aktivität, aber schnell vorbei. Lange HWZ = geringere Aktivität, aber lange Belastung.
\end{enumerate}
\end{warnbox}

\vspace{12pt}

% ═══════════════════════════════════════════════════════════════
\abschnitt{5. Benötigtes Material}

\begin{itemize}
\item Lernpfad: \texttt{lernpfad-halbwertszeit.html}
\item Arbeitsblatt: \texttt{AB-03-halbwertszeit.pdf} (Klassensatz)
\item Musterlösung: \texttt{ML-03-halbwertszeit.pdf}
\item \textbf{Würfelzerfall:} 30-40 Würfel (oder 1-2 pro Schüler)
\item PhET-Simulation: \texttt{phet.colorado.edu/de/simulations/radioactive-dating-game}
\item Millimeterpapier für Zerfallskurve
\end{itemize}

\vspace{12pt}

% ═══════════════════════════════════════════════════════════════
\abschnitt{6. Lösungen (Kurzfassung)}

\textbf{Kasten 1:} HWZ = Zeit, in der die Hälfte zerfällt. Nach 1 HWZ: 50\%. Konstant für jedes Isotop.

\textbf{Kasten 3:} Rn-220: 56s / I-131: 8 Tage / Co-60: 5,3 Jahre / Cs-137: 30 Jahre / C-14: 5730 Jahre / U-238: 4,5 Mrd. Jahre

\textbf{Übung 1:} 0→1000 (100\%) / 1→500 (50\%) / 2→250 (25\%) / 3→125 (12,5\%) / 4→62,5 (6,25\%)

\textbf{Übung 2a:} 6h = 3 HWZ → 8000→4000→2000→1000 Kerne

\textbf{Übung 2b:} 1200→600→300→150 = 3 HWZ → 3×8 = 24 Tage

\textbf{Übung 4:} 25\% = 2 HWZ → 2×5730 = 11.460 Jahre

\vspace{8pt}
\textit{→ Vollständige Lösungen siehe ML-03-halbwertszeit.pdf}

\end{document}
